\documentclass{article}
\usepackage[utf8]{inputenc}
\usepackage{amsmath}
\usepackage{amsfonts}
\usepackage[margin=0.7in]{geometry}
\usepackage{graphicx}
\usepackage{amssymb}
\usepackage{gensymb}
\usepackage{booktabs}
\usepackage{wrapfig}
\usepackage{caption}
\usepackage[export]{adjustbox}
\usepackage[pdfborder={0 0 0}]{hyperref}
\usepackage{pdfpages}
\usepackage{fancyhdr}
\usepackage{parskip}
\usepackage{lastpage}
\usepackage{listings}
\usepackage{multirow}
\usepackage{xifthen}
\usepackage{xcolor}
\usepackage{mathtools}
\graphicspath{{images/}}
\pagestyle{fancy}

\title{Økonomi Regneark}
\author{hpz230}
\date{2022-2023}
\usepackage{natbib}

\begin{document}

\maketitle

\section*{Regneark}

\section{Øvelsesark 1 November}
Partieltafled følgende funktioner: \\
Opgave 1 \\
a) Partielafled for $x_2$ \\
\begin{equation}
    u(x_1,x_2)=x_1^a*x_2^b
\end{equation}
\begin{equation}
    \frac{du}{dx_2}=b*x_1^a*x_2^{b-1}
\end{equation}
b) Partielafled for $x_1$\\ 
\begin{equation}
     u(x_1,x_2)=a*log(x_1)+b*log(x_2)
\end{equation}
\begin{equation}
    \frac{du}{dx_1}=\frac{a}{x}
\end{equation}
c) Partielafled for alle 3 variable\\
\begin{equation}
     x(m,p_1,a)=\frac{a*m}{p_1}\longrightarrow a*m*p_1^{-1}
\end{equation}
\begin{equation}
    \frac{dx}{dm}=\frac{a}{p_1}
\end{equation}
\begin{equation}
    \frac{dx}{da}=\frac{m}{p_1}
\end{equation}
\begin{equation}
    \frac{dx}{dp_1}=a*m*p_1^{-1}\longrightarrow -a*m*p_1^{-2} \longrightarrow \frac{-a*m}{p_1^2}
\end{equation}
d) Partielafled for $x_1$ \\
\begin{equation}
     y(x_1,x_2)=(x_1^a+x_2^b)^{\frac{1}{a}}
\end{equation}
\begin{equation}
     \frac{dy}{dx_1}=\frac{1}{a}(x_1^a+x_2^b)^{\frac{1-a}{a}}*a*x_1^{a-1} \longrightarrow (x_1^a+x_2^b)^{\frac{1-a}{a}}*x_1^{a-1}
\end{equation}
e) Partielafled for $x_1$
\begin{equation}
     x_2(x_1,y)=log(x_1+y)
\end{equation}
\begin{equation}
     \frac{dx_2}{dx_1}=log(x_1+y) \longrightarrow \frac{1}{x_1+y}
\end{equation}
f) Partielafled for $x$ 
\begin{equation}
     y(x,z)=f(g(x),z)
\end{equation}
\begin{equation}
     \frac{dy}{dx}=f(g(x),z)
\end{equation}
\begin{equation}
     \frac{dy}{dx}=\frac{df}{dg}f(g(x)*\frac{dg}{dx}g(x)
\end{equation}
g) Partielafled for $x_1$
\begin{equation}
     u(x_1,x_2)=min\{x_1^a,x_2\}
\end{equation}
h) Partielafled for $x_2$
\begin{equation}
     u(x_1,x_2)=x_1+x_2*log(x_2)
\end{equation}
\begin{equation}
     \frac{du}{dx_2}=1*log(x_2)+x_2*\frac{1}{x_2}\longrightarrow log(x_2)+1
\end{equation}
\\ \\
Opgave 2: \\
\begin{equation}
     u=(x_1+x_2)^a+x_1
\end{equation}
Isoler $x_2$:
\begin{equation}
     u-x_1=(x_1+x_2)^a
\end{equation}
\begin{equation}
     (u-x_1)^{\frac{1}{a}}=x_1+x_2
\end{equation}
\begin{equation}
     (u-x_1)^{\frac{1}{a}}-x_1=x_2
\end{equation}
\begin{equation}
     \frac{dx_2}{dx_1}=\frac{1}{a}*(u-x_1)^{\frac{1-a}{a}}-1
\end{equation}
Opgave 3:

\begin{equation}
    max(x,y)=a*log(x)+(1-a)*log(y)-\lambda(p_x*x+p_y*y-m)
\end{equation}

\begin{equation}
   \frac{dM}{dx}=\frac{a}{x}-\lambda*p_x
\end{equation}
\begin{equation}
   \frac{dM}{dy}=\frac{1-a}{y}-\lambda*p_y
\end{equation}
\begin{equation}
   \frac{dM}{d\lambda}=p_x*x+p_y*y-m
\end{equation}
Nu isolerer jeg lambda:
\begin{equation}
  0=\frac{a}{x}-\lambda*p_x
\end{equation}
\begin{equation}
  \lambda*p_x=\frac{a}{x}
\end{equation}
\begin{equation}
  \lambda=\frac{a}{x*p_x}
\end{equation}
Nu sætter jeg lambda ind i den y afledte funktion: \\
Jeg vil isolere y: \\
\begin{equation}
   0=\frac{1-a}{y}-\frac{a}{x*p_x}*p_y
\end{equation}
\begin{equation}
   0=\frac{1-a}{y}-\frac{a*p_y}{x*p_x}
\end{equation}
\begin{equation}
   \frac{a*p_y}{x*p_x}=\frac{1-a}{y}
\end{equation}
\begin{equation}
   \frac{1}{\frac{a*p_y}{x*p_x}}=\frac{1}{\frac{(1-a)}{y}}
\end{equation}
\begin{equation}
   \frac{x*p_x}{a*p_y}=\frac{y}{(1-a)}
\end{equation}
\begin{equation}
   \frac{x*p_x*(1-a)}{a*p_y}=y
\end{equation}
Jeg sætter y ind i funktionen for $\lambda$ igen.
\begin{equation}
   p_x*x+p_y*(\frac{x*p_x*(1-a)}{a*p_y})=m
\end{equation}
\begin{equation}
   \frac{p_x*x}{1}+\frac{x*p_x*(1-a)}{a}=m
\end{equation}
Forlænger brøken med a.
\begin{equation}
   \frac{p_x*x*a}{a}+\frac{x*p_x*(1-a)}{a}=m
\end{equation}
\begin{equation}
   \frac{p_x*x*a+x*p_x*(1-a)}{a}=m
\end{equation}




\newpage
\section*{ØvelsesOpgaver Kapitel 1 fra bogen}
1. Suppose that there were 25 people who had a reservation price of 500\$, and the 26th person had a reservation price of 200\$. What would the demand curve look like? \\ \\
25 personer, reservation price 500\$.  \\
Den 26 person, reservation price 200\$. \\
Grafen vil vise en trappe med et højt trin, hvor den starter ved 500 og går ud, så går den stejlt ned til 200 hvor den 26 person har reservationsprisen. De 25 lejer lejlighed for 500\$, og den 26 lejer for 200\$. \\ 

2. In the above example, what would the equilibrium price be if there were 24 apartments to rent? What if there were 26 apartments to rent? What if there were 25 apartments to rent? \\ \\
Ved 24 lejligheder ligger ligevægten på 500\$. Ved 25 lejligheder er der ikke nogen 25 mand til at byde 500 eller over. Den 26 mand byder 200. Derfor vil prisen faktisk ligge mellem 500\$ og 200\$. \\
26 lejligheder der ligger prisen mellem 1-200\$. \\

7. Suppose the demand curve is $D(p) = 100-2p$. What price would the monopolist set if he had 60 apartments? How many would he rent? What price would he set if he had 40 apartments? How many would he rent? \\ \\
Sætter 60 ind i funktionen og isolere p. \\
$60=100-2p$ \\
$\frac{-40}{-2}=p$ \\
$20=p$ \\ 
$60*20=1200$ \\ \\
$40=100-2p$ \\
$\frac{-60}{-2}=p$ \\
$30=p$ \\
$40*30=1200$ \\ \\
$d=100-2p$ \\ 
$d-100=-2p$ \\
$\frac{100-d}{2}=p(D)$ \\ \\
Vi skal finde profitfunktionen:
$f(D)=D*\frac{100-d}{2} = 50*D-\frac{D^2}{2}$ \\
Vi finder den afledte. \\
$f'(D)=50-d=0$ \\
$D=50$ \\
Ved at udleje 50 lejligheder, har han max profit. \\
$50=100-2p$ \\
$50-100=-2p$ \\
$\frac{50-100}{-2}=p$\\
$25=p$ \\
$50*25=1250$ \\


8. If our model of rent control allowed for unrestricted subletting, who would end up getting apartments in the inner circle? Would the outcome be Pareto efficient? \\ \\
Siden der er ubegrænset fremleje. så er det muligt at der er folk der har en lejlighed i inner-circle som hellere vil betale mindre end 500\$, og samtidigt er der andre i outer-circle der gerne vil betale 500\$. Derfor kan der indgås en handel indbyrdes. Derfor er den Pareto Inefficient. \\
\\ \\ \\ \\
1. Originally the consumer faces the budget line $p_1x_1 + p_2x_2 = m$. Then the price of good 1 doubles, the price of good 2 becomes 8 times larger, and income becomes 4 times larger. Write down an equation for the new budget line in terms of the original prices and income. \\ \\
$2*p_1x_1 + 8*p_2x_2 = 4*m$ \\
$\frac{p_1x_1}{2} + 2*p_2x_2 = m$ \\

7. If the income of the consumer increases and one of the prices decreases at the same time, will the consumer necessarily be at least as well-off? \\ \\
Indkomsten for forbrugen stiger, og en af goderne falder i pris. Budgetlinjen i dette tilfælde vil enten hælde mere eller blive fladere, og samtidigt gå lidt ud eller op ad hhv x eller y-aksen. \\


Opgave 1.3 \\
En universitetslektor lever af sherry $(S)$ og romkugler $(R)$. Prisen på sherry er
$pS = 40$, prisen på romkugler er $pR = 9$ og indkomsten er 3000 kr. pr. måned
efter øvrige udgifter er betalt. \\
1. Illustrer budgetmængden. \\ \\
$pS=40$ og $pR=9$ og $m\leq3000$ \\
$pS*40+pR*9\leq3000$ \\
$75$ cherry flasker hvis han ikke køber romkugler \\
$333.33...$ romkugler, hvis han ikke køber cherry flasker. \\



2. Lektoren kan højst drikke 1 flaske sherry pr. dag, da han ellers lugter for
meget af alkohol. Illustrer budgetmængden og udregn hvor mange romkugler
han i gennemsnit mindst spiser pr. dag. \\ \\
$40*30+pS*9=3000$ \\
Isoler Ps:\\
$pS=\frac{3000-40*30}{9}$ \\
$pS=\frac{3000-40*30}{9}$ \\
$pS\approx 200$ romkugler om måneden. \\
$329/30\approx6.5$ romkugler om dagen igennemsnit. 

\newpage
\section*{Øvelseshæfte 1.8, 1.9, 1.6 ikke g og h}
1.6)\\
a) $$ 3x_1 + 5x_2$$
Man skal partielaflede funktionen og dividere de to for at finde MRS.
$$\frac{\partial}{\partial x_1}=3$$
$$\frac{\partial}{\partial x_2}=5$$
$$MRS=\frac{3}{5}$$
b) $$ax_1 + bx_2$$
$$\frac{\partial}{\partial x_1}=a$$
$$\frac{\partial}{\partial x_2}=b$$
$$MRS=\frac{a}{b}$$
c) $$x_1^{1/3} x_2^{1/2}$$
$$\frac{\partial}{\partial x_1}=\frac{1}{3}x_1^{-2/3}x_2^{1/2}$$
$$\frac{\partial}{\partial x_2}=x_1^{1/3}\frac{1}{2}x_2^{-1/2}$$
$$MRS=\frac{\frac{1}{3}x_1^{-2/3}x_2^{1/2}}{x_1^{1/3}\frac{1}{2}x_2^{-1/2}}\longrightarrow \frac{\frac{x_2^{1/2}}{3*x_1^{2/3}}}{\frac{x_1^{1/3}}{2*x_2^{1/2}}}\longrightarrow \frac{(x_2^{1/2})*(x_1^{1/3})}{(3*x_1^{2/3})*(2*x_2^{1/2})}$$ 
d) $$x^a_1 x^b_2$$
$$\frac{\partial}{\partial x_1}=ax_1^{a-1}x_2^b$$
$$\frac{\partial}{\partial x_2}=x_1^{a}bx_2^{b-1}$$
e) $$(x_1^{1/3}+x_2^{1/3})^3$$
$$\frac{\partial}{\partial x_1}=3(x_1^{1/3}+x_2^{1/3})^2*\frac{1}{3}x_1^{-2/3}$$
$$\frac{\partial}{\partial x_2}=3(x_1^{1/3}+x_2^{1/3})^2*\frac{1}{3}x_2^{-2/3}$$
f) $$(x_1^a+x_2^a)^{1/a}$$
$$\frac{\partial}{\partial x_1}=\frac{1}{a}(x_1^a+x_2^a)^{-1/a}*ax_1^{a-1}$$
$$\frac{\partial}{\partial x_1}=\frac{1}{a}(x_1^a+x_2^a)^{-1/a}*ax_2^{a-1}$$
g) $$-e^{-x_1x_2}$$
$$\frac{\partial}{\partial x_1}=x_2e^{-x_1x_2}$$
$$\frac{\partial}{\partial x_2}=x_1e^{-x_1x_2}$$
\newpage
1.8)\\
1. Opstil Elviras problem og den tilhørende Lagrange funktion og udregn de tilhørende første-ordens betingelser.\\ \\
Vi ved at $U=x_1*x_2$ Elviras problem er: at vi ønsker at finde hendes max, med hensyn til $x_1x_2$ under hendes budgetbetingelse. 
$$x_2=\frac{U}{x_1}\longrightarrow U*x_1^{-1}$$
$$\frac{dx_2}{dx_1}=U*(-1)*x_1^{-2}$$ 
$$\frac{d^2x_2}{dx_1}=U*(-1)*(-2)*x_1^{-3}\longrightarrow k*2*x_1^{-3}$$
Den dobbeltafledte er positiv, så vi kan benytte Lagrange.\\


2. Find Elviras efterspørgselsfunktioner, start evt. med $p_1 = p_2 = m = 1$.\\ \\
$$L(x_1,x_2,\lambda)=x_1*x_2-\lambda(p_1x_1+p_2x_2-m)$$
$$\frac{\partial L}{\partial x_1}=x_2-\lambda*p_1$$
$$\frac{\partial L}{\partial x_2}=x_1-\lambda*p_2$$
$$\frac{\partial L}{\partial \lambda}=p_1x_1+p_2x_2-m$$
\\ \\
$$x_2-\lambda*p_1\longrightarrow\frac{x_2}{p_1}=\lambda$$
$$x_1-\frac{x_2}{p_1}*p_2\longrightarrow x_1=\frac{x_2*p_2}{p_1}$$ \\ \\
$$p_1*\frac{x_2*p_2}{p_1}+p_2x_2=m\longrightarrow x_2*p_2+x_2*p_2\longrightarrow 2*x_2*p_2=m\longrightarrow x_2=\frac{m}{2*p_2}$$
\\\\
$$p_1x_1+p_2\frac{m}{2*p_2}=m\longrightarrow p_1x_1+\frac{m}{2}=m\longrightarrow p_1x_1=m-\frac{m}{2}\longrightarrow x_1=\frac{m}{2*p_1}$$ \\ \\

$$x_1(,p_1,p_2,m)=\frac{m}{2*p_1}$$


\newpage
3. Tegn Engelkurven for vare 1 - er vare 1 inferiør?\\ \\
$$x_1(m)=m\frac{1}{p_1*2}$$
Den har punkt i $(1,\frac{1}{p_1*2})$ Engelkurven går fra origo og ud mod punktet.



4. Tegn efterspørgselskurven for vare 1 - er vare 1 et Giffen gode?\\ \\
Efterspørgselskurven for $x_1$: \\
$$x_1(p_1)=\frac{m}{p_1*2}=\frac{m}{2}*\frac{1}{p_1}$$
Efterspørgselskurven laver en konveks differencekurve, som går fra y mod x.

$$\frac{dx_1}{dp_1}=\frac{m}{2}*(-1)*p_1^{-2}=\frac{-m}{2*p_2^2}$$
Derfor er det ikke en giffengode, men en ordinær gode, fordi det er mindre end 0.
\end{document}
